\documentclass[10pt,a4paper]{article}
\usepackage{amsmath,amssymb,amsthm,graphicx}
% Minimal retained preamble for lane-p2 cards taken from
% sources/ahl-16-fixed-points-2019/editorial-preamble.tex
% (doi:10.5802/ahl.16, Annales Henri Lebesgue 2 (2019) 149-185, CC BY 4.0).
%
% Copied verbatim from the editorial preamble: the package list, the
% theorem declarations under the shared counter and the macro block, so that
% every macro and environment of the retained spans is available.
%
% Omitted and why: the class line (the assembler supplies the class); the
% a4paper geometry line and the fontspec block with its explicit lmroman10
% files (encodings and fonts fight the forced xelatex engine, and the
% assembler supplies the page setup).
\usepackage{amsmath,amssymb,amsthm,mathtools,graphicx,xcolor,hyperref,xurl,xspace,ifthen}
\newtheorem{theo}{Theorem}[section]
\newtheorem{lemm}[theo]{Lemma}
\newtheorem{theorem}[theo]{Theorem}
\newtheorem*{remas*}{Remarks}
\providecommand{\og}{«\nobreakspace}
\providecommand{\fg}{\nobreakspace»}
\newtheorem{lem}[theo]{Lemma}
\newtheorem{lemma}[theo]{Lemma}
\newtheorem{prop}[theo]{Proposition}
\newtheorem{coro}[theo]{Corollary}
\newtheorem{corol}[theo]{Corollary}
\newtheorem{conj}[theo]{Conjecture}
\newtheorem{Proposition}[theo]{Proposition}
\newtheorem{theoreme}[theo]{Theorem}
\theoremstyle{definition}
\newtheorem{defi}[theo]{Definition}
\newtheorem{exam}[theo]{Example}
\newtheorem{exem}[theo]{Example}
\newtheorem{rema}[theo]{Remark}
\newtheorem{nota}[theo]{Notation}
\newtheorem*{rema*}{Remark}
\newtheorem*{defi*}{Definition}
\newtheorem*{theo*}{Theorem}
\newtheorem*{lemm*}{Lemma}
\newtheorem*{prop*}{Proposition}
\newcommand{\enoncename}{}
\newtheorem{corpusenonce}[theo]{\enoncename}
\newenvironment{enonce}[2][]{\renewcommand{\enoncename}{#2}\begin{corpusenonce}}{\end{corpusenonce}}
\newenvironment{enonce*}[2][]{\par\medskip\noindent\textbf{#2.}\itshape}{\par\medskip}
\newenvironment{proofc}[1][\proofname]{\begin{proof}[#1]}{\end{proof}}
\newcommand{\Mc}{\mathcal{M}}
\newcommand{\pointir}{.\ ---\ }
\providecommand{\mathof}[1]{\mathrm{#1}}

\numberwithin{equation}{section}
\numberwithin{figure}{section}
\numberwithin{table}{section}
\def\endappendix{}
%~Mouliné par MaN_auto v.0.21.4 2019-04-30 09:02:09


\usepackage[all]{xy}

\hyphenation{auto-morphism
every
Roelcke
peri-od-ic
peri-od-icy
rep-re-sen-ta-tion
}








\newcommand{\A}{\mathcal A}
\newcommand{\N}{\mathbb{N}}
\newcommand{\Z}{\mathbb{Z}}
\newcommand{\Q}{\mathbb{Q}}
\newcommand{\K}{\mathcal{K}}
\newcommand{\CC}{\mathbb{C}}

\newcommand*{\mA}{\mathbf{A}}
\newcommand*{\mB}{\mathbf{B}}
\newcommand*{\mC}{\mathbf{C}}
\newcommand*{\mD}{\mathbf{D}}
\newcommand*{\mE}{\mathbf{E}}
\newcommand*{\mF}{\mathbf{F}}
\newcommand*{\mH}{\mathbf{H}}
\newcommand*{\mZ}{\mathbf{Z}}
%\newcommand*{\mA}{\mathbf{A}}
%\newcommand*{\mA}{\mathbf{A}}

\newcommand{\restrict}[2]{#1 \upharpoonright #2}

\newcommand{\e}{\varepsilon}

\DeclareMathOperator{\Aut}{Aut}
\DeclareMathOperator{\Stab}{Stab}
\DeclareMathOperator{\Age}{Age}

\newcommand{\actson}{\curvearrowright}

\newcommand{\RUCB}{\mathrm{RUC}_{b}}
\newcommand{\ROB}{\mathrm{Ro}_{b}}
\newcommand{\WAP}{\mathrm{WAP}}
\newcommand{\AP}{\mathrm{AP}}

\graphicspath{{./figures/}}

\newcommand*{\mk}{\mkern -1mu}
\newcommand*{\Mk}{\mkern -2mu}
\newcommand*{\mK}{\mkern 1mu}
\newcommand*{\MK}{\mkern 2mu}

\hypersetup{urlcolor=purple, linkcolor=blue, citecolor=red}

\newcommand*{\romanenumi}{\renewcommand*{\theenumi}{\roman{enumi}}}
\newcommand*{\Romanenumi}{\renewcommand*{\theenumi}{\Roman{enumi}}}
\newcommand*{\alphenumi}{\renewcommand*{\theenumi}{\alph{enumi}}}
\newcommand*{\Alphenumi}{\renewcommand*{\theenumi}{\Alph{enumi}}}
\let\oldexists\exists
\renewcommand*{\exists}{\mathrel{\oldexists}}
\let\oldforall\forall
\renewcommand*{\forall}{\mathrel{\oldforall}}


\usepackage{cleveref}
\crefname{originalproposition}{original proposition}{original propositions}
\Crefname{originalproposition}{Original proposition}{Original propositions}
\makeatletter
\newlabel{cor:fp}{{2.3}{source}{}{sourcepointer-1}{}}
\newlabel{cor:fp@cref}{{[originalproposition][0][]2.3}{[1][0][]source}{}{}{}}
\makeatother
\crefname{originalproposition}{original proposition}{original propositions}
\Crefname{originalproposition}{Original proposition}{Original propositions}
\makeatletter
\newlabel{prop:OSRamsey}{{3.5}{source}{}{sourcepointer-2}{}}
\newlabel{prop:OSRamsey@cref}{{[originalproposition][0][]3.5}{[1][0][]source}{}{}{}}
\makeatother
\title{A fixed point for the flow on the spectrum of a left-invariant algebra is equivalent to the Ramsey property for its colorings}
\author{Lionel Nguyen Van Thé}
\date{doi:10.5802/ahl.16; CC BY 4.0}
\begin{document}
\maketitle
\noindent\textit{Editorial note.} The counted unit is the article's own Proposition cor:fpRP at source lines 378-384 together with its proof at source lines 386-388, a straight adjacent pairing in the source file. The source is the publisher version of the article in Annales Henri Lebesgue 2 (2019), pages 149-185, DOI 10.5802/ahl.16, released under CC BY 4.0.\par\medskip
\noindent\textit{Editorial note.} The statement environment carries the label cor:fpRP but the article prints it as a proposition, not as a corollary; the environment kind used for the pointer bookkeeping is therefore a proposition, which is what the reader sees. The same environment holds an enumerated list whose items are lettered by the article's own numbering, which is why the retained statement spans the whole environment rather than a single line.\par\medskip
\noindent\textit{Editorial note.} The licence evidence used here is the observation file kept with the source, license-observation.json, which records the licence as CC BY 4.0, points at the inherited publisher-version evidence and states that the landing page displays no licence line; the underlying evidence files are the same ones the earlier cards of this source rely on. No acceptance record is written anywhere.\par\medskip
\noindent\textit{Editorial note.} Both pointers of this card were numbered the way the article numbers its statements: its theorem-like environments share one counter attached to the section, so the printed number is the section together with the place of the statement inside it. The numbers were read off an enumeration of every numbered theorem-like environment of the source in order, which gives 2.3 for the fixed point proposition and 3.5 for the oscillation stability proposition. An earlier group of cards of this source was re-assembled with the same printed numbering before this card was written; the method and the full table are recorded in this lane's STATUS. The anchors recorded point at each statement's own lines.\par\medskip
\noindent\textit{Editorial note.} The gate raised one manual-review note: the next proof environment begins at line 390. That block proves the article's next theorem and is not part of this card; the proof retained here is closed at line 388. Read from the source before writing, and no further action was needed.\par\medskip
\noindent\textit{Editorial note.} Neither the retained statement nor the retained proof contains a citation, so this card needs no bibliography entry.\par\medskip
\noindent\textit{Editorial note.} The census package marked this candidate free of missing citation keys and listed the statement span as 378-384 and the proof span as 386-388; both were re-read from the source before the card was written.\par\medskip
\noindent\textit{Editorial note.} rqual sv-mismatch batch fix 2026-10-14 (Chief ruling: re-assemble unpublished only): source\_version aligned verbatim to the source dir license-observation.json sv (package evidence authority); revision bumped for re-assembly. Recipe from lane-s3.\par\medskip
\noindent\textit{Allowed standard prerequisites (editorial).}\begin{itemize}
\item Left-invariant C-star subalgebras of the bounded right-uniformly continuous functions and their spectra. For a topological group the article considers the unital, left-invariant, closed C-star subalgebras of the algebra of bounded right-uniformly continuous functions, the flow of the group on the spectrum of such an algebra, and the finite colorings that belong to the algebra, so that a finite coloring is exactly the kind of element whose topological behaviour the flow records. The author's own definitions are the ones used by the counted statement and proof.. The counted proposition is the general statement that ties a fixed point of this flow to the Ramsey property of the structure for the colorings that lie in the algebra.
\item Finite oscillation stability and its equivalent Ramsey form. A finite family of functions is finitely oscillation stable when for every finite set of group elements and every precision there is a point at which all the functions in the family vary by less than the precision; the article's oscillation stability proposition states that for a finite set of finite colorings this is equivalent to finding, for every finite substructure, a copy of it on which every coloring of the family is constant. The author's own definitions and the surrounding results are the ones used by the counted statement and proof.. The proof uses the equivalent forms in both directions: the fixed point is first translated into finite oscillation stability, and the combinatorial form of that property is then recognised as the Ramsey property.
\end{itemize}\medskip
\ifdefined\hypertarget\hypertarget{sourcepointer-1}{}\fi
\noindent\textit{Source reference.} Original proposition 2.3: sources/ahl-16-fixed-points-2019/body-01.tex, source lines 260--266. Navigation pointer invoked by the retained proof for the translation of the fixed point of the flow into finite oscillation stability. The statement carries the label of a corollary but is printed as a proposition by the article; it is an earlier result and is not retained in this card, so the reference is declared as a pointer carrying the number the article itself prints for that statement, that is the section and the place of the statement inside that section, counted by enumerating every numbered theorem-like environment of the source in order: it is 2.3; the cited statement is an allowed input and none of it is reproduced inside the span.\par\medskip
\ifdefined\hypertarget\hypertarget{sourcepointer-2}{}\fi
\noindent\textit{Source reference.} Original proposition 3.5: sources/ahl-16-fixed-points-2019/body-01.tex, source lines 368--374. Navigation pointer invoked by the retained proof to recognise finite oscillation stability of a finite set of finite colorings as the combinatorial Ramsey condition. That proposition is an earlier result of the article and is not retained in this card, so the reference is declared as a pointer carrying the number the article itself prints for that statement, that is the section and the place of the statement inside that section, counted by enumerating every numbered theorem-like environment of the source in order: it is 3.5; the cited statement is an allowed input and none of it is reproduced inside the span.\par\medskip
\begin{prop}\label{cor:fpRP}
Let $\mF$ be a Fra\"{\i}ss\'e structure, and let $\A$ be a unital, left-invariant, closed C$^{*}$-subalgebra of $\RUCB(\Aut(\mF))$. Assume that finite colorings are dense in $\A$. TFAE:
\begin{enumerate}
\item \label{3.7.1} The flow $\Aut(\mF)\actson \Aut(\mF)^{\A}$ has a fixed point.
\item \label{3.7.2} The structure $\mF$ has the Ramsey property for colorings in $\A$.
\end{enumerate}
\end{prop}

\begin{proof}
Thanks to Proposition~\ref{cor:fp}, the flow $\Aut(\mF)\actson \Aut(\mF)^{\A}$ has a fixed point iff every finite $\mathcal F\subset \A$ is finitely oscillation stable. Because finite colorings are dense in $\A$, this holds iff every finite set $\mathcal C \subset \A$ of finite colorings is oscillation stable. This is equivalent to $\mF$ having the Ramsey property for colorings in $\A$ thanks to Proposition~\ref{prop:OSRamsey}.
\end{proof}

\end{document}
